% !TEX program = pdflatex
% Kompiliert mit: pdflatex (2x)  | optional: biber (falls biblatex genutzt wird)

\documentclass[
  12pt,
  a4paper,
  oneside,
  headings=big,
  parskip=half,
  bibliography=totoc
]{scrreprt}

% --- Sprache & Typografie ---
\usepackage[T1]{fontenc}
\usepackage[utf8]{inputenc} % ok für pdflatex
\usepackage[ngerman]{babel}
\usepackage{microtype}
\usepackage{csquotes}

% --- Layout ---
\usepackage[a4paper,margin=2.5cm]{geometry}
\usepackage{setspace}
\onehalfspacing

% --- Mathe, Tabellen, Grafiken ---
\usepackage{amsmath,amssymb}
\usepackage{graphicx}
\usepackage{booktabs}
\usepackage{array}
\usepackage{longtable}
\usepackage{tabularx}

% --- Listen & Links ---
\usepackage{enumitem}
\setlist{nosep}
\usepackage[hidelinks]{hyperref}
\usepackage{url}

% --- KOMA: bessere Überschriftenabstände (optional) ---
% \RedeclareSectionCommand[beforeskip=0.8\baselineskip,afterskip=0.6\baselineskip]{chapter}

% --- Optionale Literaturverwaltung (wenn du später BibLaTeX nutzen willst) ---
% \usepackage[backend=biber,style=authoryear,maxcitenames=2]{biblatex}
% \addbibresource{references.bib}

% --- Metadaten / Titelseite ---
\newcommand{\thesistitle}{Risikomanagement in Data-Science-Projekten}
\newcommand{\thesissubtitle}{Einordnung, Methoden und Strategieentwicklung entlang des DASC-PM}
\newcommand{\thesisauthor}{\textless Dein Name\textgreater}
\newcommand{\thesisprogram}{Masterstudium \textless Studiengang\textgreater}
\newcommand{\thesiscourse}{Modul: Projektmanagement in Data-Science-Projekten}
\newcommand{\thesissupervisor}{\textless Betreuer:in\textgreater}
\newcommand{\thesisdate}{15. April 2026}

\title{\thesistitle\\\large \thesissubtitle}
\author{\thesisauthor}
\date{\thesisdate}

\begin{document}

% --- Titelseite (einfach, gut editierbar) ---
\begin{titlepage}
  \centering
  \vspace*{2cm}

  {\LARGE\bfseries \thesistitle\par}
  \vspace{0.5cm}
  {\Large \thesissubtitle\par}

  \vspace{1.5cm}
  {\large \thesisauthor\par}

  \vfill

  \begin{tabular}{@{}ll@{}}
    Studiengang: & \thesisprogram\\
    Veranstaltung: & \thesiscourse\\
    Betreuung: & \thesissupervisor\\
    Datum: & \thesisdate\\
  \end{tabular}

  \vspace*{2cm}
\end{titlepage}

% --- Optional: Abstract (de) ---
\chapter*{Abstract}
\addcontentsline{toc}{chapter}{Abstract}
% Kurzfassung in 150--250 Wörtern: Problem, Vorgehen, Ergebnis, Beitrag.
\noindent\textless Hier kommt dein Abstract hin.\textgreater

\tableofcontents
\clearpage

% --- Optional: Abkürzungsverzeichnis (minimal, ohne zusätzliche Pakete) ---
\chapter*{Abkürzungsverzeichnis}
\addcontentsline{toc}{chapter}{Abkürzungsverzeichnis}
\begin{description}[leftmargin=3.2cm,style=nextline]
  \item[DASC-PM] \textless Bedeutung eintragen\textgreater
  \item[DS] Data Science
  \item[PM] Projektmanagement
  \item[RM] Risikomanagement
\end{description}
\clearpage

% ==========================================================
% Gliederung (Baukasten)
% ==========================================================

\chapter{Einleitung}
\section{Motivation und Problemstellung}
\textless Text\textgreater

\section{Zielsetzung und Forschungsfragen}
\textless Text\textgreater

\section{Vorgehen und Aufbau der Arbeit}
\textless Text\textgreater


\chapter{Grundlagen des Risikomanagements}
\section{Begriffe und Definitionen}
\textless Text\textgreater

\section{Risikotypen in Data-Science-Projekten}
\textless Text\textgreater

\section{Rollen, Verantwortlichkeiten und Governance}
\textless Text\textgreater


\chapter{Methoden des Risikomanagements}
\section{Überblick: Prozessschritte und Artefakte}
\textless Text\textgreater

\section{Methode 1: Risikobewertung}
\subsection{Identifikation}
\textless Text\textgreater

\subsection{Analyse und Bewertung (qualitativ/quantitativ)}
\textless Text\textgreater

\subsection{Beispiel-Artefakt: Risk Register (Template)}
\begin{longtable}{@{}p{3.0cm}p{3.2cm}p{2.1cm}p{2.1cm}p{4.2cm}@{}}
\toprule
\textbf{Risiko} & \textbf{Ursache/\allowbreak Trigger} & \textbf{Eintritt} & \textbf{Auswirkung} & \textbf{Maßnahmen /\allowbreak Owner} \\
\midrule
\endfirsthead
\toprule
\textbf{Risiko} & \textbf{Ursache/\allowbreak Trigger} & \textbf{Eintritt} & \textbf{Auswirkung} & \textbf{Maßnahmen /\allowbreak Owner} \\
\midrule
\endhead
\textless Risiko 1\textgreater & \textless Trigger\textgreater & H/M/L & H/M/L & \textless Maßnahme + Verantwortliche:r\textgreater \\
\textless Risiko 2\textgreater & \textless Trigger\textgreater & H/M/L & H/M/L & \textless Maßnahme + Verantwortliche:r\textgreater \\
\bottomrule
\end{longtable}

\section{Methode 2: Risikominimierungsstrategien}
\subsection{Vermeiden, Vermindern, Übertragen, Akzeptieren}
\textless Text\textgreater

\subsection{Kontrollen, Monitoring und Eskalation}
\textless Text\textgreater


\chapter{Strategieentwicklung entlang des DASC-PM}
\section{Annahmen: Phasen/Artefakte des DASC-PM}
\textless Text\textgreater

\section{Risikomanagement-Integration je Phase}
\textless Text\textgreater

\section{Vorschlag: RM-Playbook als wiederverwendbarer Baukasten}
\textless Text\textgreater


\chapter{Case-Study}
\section{Auswahl, Beschreibung und Ziel der Case-Study}
\textless Text\textgreater

\section{Anwendung der Strategie auf die Case-Study}
\textless Text\textgreater

\section{Case-Study 1: Anwendung und Bewertung}
\textless Text\textgreater

\section{Case-Study 2: Anwendung und Bewertung}
\textless Text\textgreater

\section{Case-Study 3: Anwendung und Bewertung}
\textless Text\textgreater


\chapter{Diskussion}
\section{Kann man das DASC-PM um ein Risikomanagement erweitern oder sollte man ein parallel laufendes Modell nutzen?}
\textless Text\textgreater

\section{Limitationen und Validität der Ergebnisse}
\textless Text\textgreater


\chapter{Fazit}
\section{Zusammenfassung der Ergebnisse}
\textless Text\textgreater

\section{Implikationen für Data-Science-Projektmanagement}
\textless Text\textgreater

\section{Ausblick}
\textless Text\textgreater


% ==========================================================
% Literatur (self-contained) – später gern auf BibLaTeX umstellen
% ==========================================================
\begin{thebibliography}{9}
  \bibitem{placeholder} \textless Autor\textgreater. \textit{\textless Titel\textgreater}. \textless Verlag/Journal\textgreater, \textless Jahr\textgreater.
\end{thebibliography}

% --- Wenn du BibLaTeX aktivierst, ersetze oben durch: ---
% \printbibliography


% ==========================================================
% Anhang (optional)
% ==========================================================
% \appendix
% \chapter{Zusätzliche Materialien}
% \textless Text\textgreater

\end{document}
